\documentclass[11pt]{article}
\usepackage[T1]{fontenc}
\usepackage{lmodern,amsmath,amssymb,amsthm,mathtools,microtype}
\usepackage[margin=1.02in]{geometry}
\usepackage[hidelinks]{hyperref}
\hypersetup{pdftitle={Unimodality of the q-Fibonomial coefficients},pdfauthor={Kalle Nikula-Gill}}
\usepackage{enumitem}
\newtheorem{theorem}{Theorem}
\newtheorem{lemma}[theorem]{Lemma}
\newcommand{\Z}{\mathbb Z}
\newcommand{\N}{\mathbb N}
\newcommand{\dd}{\Delta}
\newcommand{\abs}[1]{\lvert#1\rvert}
\DeclareMathOperator{\ord}{ord}
\title{Unimodality of the \(q\)-Fibonomial coefficients}
\author{Kalle Nikula-Gill}
\date{}
\begin{document}
\maketitle
\begin{abstract}
We give a computer-assisted proof of the unimodality conjecture of
Bergeron, Ceballos and K\"ustner. A recurrence reduces the inductive
step to \(F_m\) consecutive differences at the midpoint. The last three
Fibonacci factors have pairwise coprime denominator lengths, so the
difference is an explicit positive linear term plus simple
root-of-unity contributions. We bound the latter uniformly for
\(m\ge40\). The remaining arithmetic consists of 1,424 integer
inequalities and 28 exact coefficient computations.
\end{abstract}

\section{Statement and polynomial model}
Let \(F_0=0,\ F_1=1,\ F_{j+2}=F_{j+1}+F_j\). Write
\([a]_q=\sum_{j=0}^{a-1}q^j\) and \(H_t=\prod_{j=1}^{t}[F_j]_q\),
with \(H_0=1\). The \(q\)-Fibonomial polynomial is
\begin{equation}\label{eq:def}
 P_{m,n}(q)=\frac{H_{m+n}}{H_mH_n}
 =\prod_{j=1}^{m}\frac{1-q^{F_{n+j}}}{1-q^{F_j}} .
\end{equation}
Bergeron, Ceballos and K\"ustner conjectured unimodality
\cite[Conjecture~2.5]{BCK}. This is Conjecture~1.1 of Connelly, Ito,
Martinez, Shevchenko and Yang, who proved the cases
\(\min(m,n)\le3\) \cite[Theorem~1.2]{CIMSY}.

\begin{theorem}\label{thm:main}
For every \(m,n\ge0\), write \(P_{m,n}(q)=\sum_{i=0}^d c_iq^i\).
With \(h=\lfloor d/2\rfloor\), the coefficients satisfy
\(c_0\le\cdots\le c_h\ge c_{h+1}\ge\cdots\ge c_d\).
\end{theorem}

We record why \eqref{eq:def} is a polynomial. Consecutive Fibonacci
numbers are coprime by the Euclidean algorithm, and induction gives
\begin{equation}\label{eq:add}
 F_{a+t}=F_{a+1}F_t+F_aF_{t-1}\qquad(t\ge1).
\end{equation}
If \(e>1\) divides some positive-index Fibonacci number, let \(r\) be
its least positive index of divisibility. Then \(e\mid F_j\) exactly
when \(r\mid j\). In fact \(F_{r+t}\equiv F_{r+1}F_t\bmod e\), also
at \(t=0\), and \(F_{r+1}\) is invertible modulo \(e\). Reduction of
\(j\) modulo \(r\) proves the assertion. Each root of \(1-q^a\) is
simple for \(a>0\). At a root of order \(e>1\) of the denominator of
\eqref{eq:def}, the interval \([1,m]\) contains \(\lfloor m/r\rfloor\)
multiples of \(r\), and \([n+1,n+m]\) contains at least as many. At
\(1\) there are \(m\) factors on each side. This proves divisibility
over \(\mathbb C\). Expansion at zero gives integer coefficients and
constant term~1.

The factorial formula is symmetric in \(m,n\). Each \([a]_q\) is
reciprocal of degree \(a-1\), so \(c_i=c_{d-i}\). Summing the
recurrence gives \(\sum_{j=1}^t F_j=F_{t+2}-1\), hence
\begin{equation}\label{eq:degree}
 d=\sum_{j=1}^m(F_{n+j}-F_j)
   =F_{m+n+2}-F_{m+2}-F_{n+2}+1.
\end{equation}
In particular \(P_{0,n}=P_{n,0}=1\). Extend all coefficient sequences
by zero to indices in \(\Z\) and put
\(\dd P(i)=[q^i](1-q)P=c_i-c_{i-1}\). Reciprocity gives
\(\dd P(d+1-i)=-\dd P(i)\), including the middle zero when \(d\) is
odd. It suffices to prove \(\dd P(i)\ge0\) for \(0\le i\le h\).

\section{Reduction to the midpoint}
Fix \(n\) and induct on \(m\le n\). For \(m=1\) we have \([F_{n+1}]_q\).
For \(m=2\) we have \([a]_q[b]_q\), where \(a=F_{n+1}\le b=F_{n+2}\).
Its differences through the midpoint are nonnegative, since
\((1-q)[a]_q[b]_q=(1-q^b)[a]_q\) and
\(\lfloor(a+b-2)/2\rfloor<b\).
For \(m=3\), put \(C=F_{n+1}, B=F_{n+2}, A=F_{n+3}=B+C\).
Here \(h=A-2\). The series \(1/((1-q)(1-q^2))\) has coefficients
\(t(i)=\lfloor i/2\rfloor+1\) for \(i\ge0\) and \(t(i)=0\)
otherwise. Through \(h\) the differences are
\(t(i)-t(i-C)-t(i-B)\). For \(i<B\) use monotonicity of \(t\).
For \(B\le i\le A-2\), use
\[
 \left\lfloor\tfrac{i-C}{2}\right\rfloor+
 \left\lfloor\tfrac{i-B}{2}\right\rfloor
 \le\left\lfloor\tfrac{2i-A}{2}\right\rfloor
 \le\left\lfloor\tfrac i2\right\rfloor-1 .
\]

For \(m\ge4\), write \(f=P_{m-1,n},\ g=P_{m,n},\
D=\deg f,\ A=F_{m+n},\ w=F_m\). Then \(d=D+A-w\), and
\begin{equation}\label{eq:diff}
 \dd g(i)-\dd g(i-w)=\dd f(i)-\dd f(i-A).
\end{equation}
Indeed \((1-q^w)g=(1-q^A)f\). By \eqref{eq:degree},
\(D<F_{m+n+1}<2A\). For \(0\le i\le D/2\) we have \(i<A\),
so the right side of \eqref{eq:diff} is nonnegative by induction.
For \(D/2<i\le d/2\) we have \(\dd f(i)\le0\) by reciprocity and
\(i-A<D/2\), so the right side is nonpositive. Along each residue
class modulo \(w\), the values of \(\dd g\) therefore rise from zero
and then fall through \(h\). The first part is nonnegative, and the
second part is bounded below by its last value. It suffices to check
\begin{equation}\label{eq:window}
 h-w+1\le i\le h .
\end{equation}

\section{The central difference}
Put \(C=F_{m+n-2},\ B=F_{m+n-1},\ A=B+C\), and
\(u=F_{m-2},\ v=F_{m-1},\ w=u+v\). Write \(Q=P_{m-3,n}\) and
\(T=\deg Q=B-F_{n+2}-v+1\). Thus \(d=T+2A-2w\).
Every \(i\) in \eqref{eq:window} satisfies
\begin{equation}\label{eq:indices}
 B+T\le i<A+C .
\end{equation}
For the upper bound use \(T<B<2C\). For the lower bound,
\(m,n\ge4\) and \(n\ge m\) give
\[
 2C-T=F_{m+n-4}+F_{n+2}+v-1
 \ge w+(2w+v)+v-1\ge4w,
\]
since \(2v-w=F_{m-3}\ge1\). Substitute
\(h=A-w+\lfloor T/2\rfloor\) in \eqref{eq:window}.

Define \(G(q)=(1-q)/((1-q^u)(1-q^v)(1-q^w))\). The three top factors
in the numerator multiply to
\(1-q^C-q^B+q^{A+C}+q^{A+B}-q^{2A}\).
By \eqref{eq:indices}, only \(1-q^C-q^B\) contributes to \(\dd g(i)\).
Further, \(i-j,\ i-B-j,\ i-C-j\) are nonnegative for every
\(0\le j\le T\). This permits the use of the same partial fraction
coefficient formula in all three terms.

The lengths \(u,v,w\) are pairwise coprime. The proper rational
function \(G\) has a double pole at 1 and simple poles at the
nontrivial roots of its denominator. Writing \(z=1-q\), we obtain
from \(1-q^U=Uz-U(U-1)z^2/2+O(z^3)\) the principal part
\[
 \frac1{uvw}\left(\frac1{(1-q)^2}
       +\frac{2w-3}{2(1-q)}\right).
\]
Its coefficient at \(t\ge0\) is \((t+(2w-1)/2)/(uvw)\).
Since \(Q\) is reciprocal, \(\sum_{j=0}^T j[q^j]Q=TQ(1)/2\).
Convolution with \((1-q^C-q^B)Q\) now gives
\begin{equation}\label{eq:main}
 \dd g(i)=\frac{Q(1)}{uvw}\left(\frac{d+1}{2}-i\right)+E_i .
\end{equation}
Indeed the linear expression inside the convolution at \(j\) becomes
\(A-i+j-(2w-1)/2\). As \(i\le h\), the displayed main term is at
least \(Q(1)/(2uvw)\), and \(Q(1)>0\) by \eqref{eq:def}.

Let \(M\) be the largest \(\abs{Q(\zeta)}/Q(1)\) over the nontrivial
roots in question. For each cyclic permutation \((U,V,W)\) of
\((u,v,w)\) and each \(\zeta\ne1\) with \(\zeta^U=1\), the
corresponding summand in \(E_i\) is
\[
 \frac{(1-\zeta)(1-\zeta^C-\zeta^B)Q(\zeta)\zeta^{-i}}
      {U(1-\zeta^V)(1-\zeta^W)}.
\]
This follows by taking the limit of \((1-q/\zeta)G(q)\) at \(\zeta\)
and expanding \(1/(1-q/\zeta)\) at zero. There is no polynomial
part in the partial fraction expansion. The first two numerator
factors have product of absolute value at most \(6\).
As \(V\equiv\pm W\bmod U\) and \(\gcd(U,V)=1\), the sum over these
roots is bounded using
\[
 \sum_{\zeta^U=1,\,\zeta\ne1}\frac1{\abs{1-\zeta}^2}
 =\frac{U^2-1}{12}.
\]
For completeness let \(H(x)=1+x+\cdots+x^{U-1}\). With
\(a_\zeta=1/(1-\zeta)\) we have
\(\bar a_\zeta=1-a_\zeta\).
Factoring \(H\) gives \(\sum a_\zeta=H'(1)/H(1)=(U-1)/2\) and
\(\sum a_\zeta^2=(\sum a_\zeta)^2-H''(1)/H(1)\), where
\(H''(1)/H(1)=(U-1)(U-2)/3\). Subtraction gives the identity, also
when \(U=1\). Summing over the three lengths yields
\[
 \abs{E_i}\le M Q(1)\sum_{U=u,v,w}\frac{U^2-1}{2U}
 \le wM Q(1).
\]
Here the sum is over the three positions. By \eqref{eq:main}, strict
positivity on the window follows if
\begin{equation}\label{eq:target} 2uvw^2 M<1. \end{equation}

\section{A bound at each root}
Let \(\zeta\) have order \(e>1\) dividing \(u\), \(v\), or \(w\).
Write \(p=m-3\) and
\[
 s=\#\{1\le j\le p:e\nmid F_j\},\qquad
 k=\#\{2\le j\le p:2F_j<e\}.
\]
We claim that
\begin{equation}\label{eq:root}
 \frac{\abs{Q(\zeta)}}{Q(1)}
 \le \frac1{k!}\left(\frac{e}{2F_{n+1}}\right)^s .
\end{equation}
Use the rank \(r\) of divisibility from Section~1. In \eqref{eq:def}
for \(Q\), vanishing factors in the denominator occur at the multiples
of \(r\) in \([1,p]\). The numerator has at least as many vanishing
factors. If it has more then \(Q(\zeta)=0\). In case of equality,
the \(\ell\)-th deleted numerator position \(c+(\ell-1)r\), for
\(1\le c\le r\), is at most the \(\ell\)-th deleted denominator
position \(\ell r\). In every prefix the numerator thus has at least
as many deletions. Its surviving indices \(i\) are correspondingly
at least the denominator indices \(j\) in increasing order.

Normalize each factor \(1-q^a\) by dividing it by \(a\). This
computes \(Q(q)/Q(1)\). At a zero the normalized derivative is
\(-\zeta^{a-1}=-\zeta^{-1}\), so each deleted pair contributes
absolute ratio one. For each surviving pair,
\(F_{n+i}\ge F_{n+j}\ge F_{n+1}F_j\) by \eqref{eq:add}.
Write \(\zeta=\exp(2\pi\mathrm{i}b/e)\) with \(\gcd(b,e)=1\), and
let \(r_j\) be the distance from \(bF_j\) to the nearest multiple
of \(e\). Then \(1\le r_j\le e/2\) and
\[
 \abs{1-\zeta^{F_j}}=2\sin(\pi r_j/e)\ge 4r_j/e
\]
by concavity of sine on \([0,\pi/2]\). The numerator factor has
absolute value at most 2 before normalization. Each pair therefore
contributes at most \(e/(2F_{n+1}r_j)\).
The \(k\) Fibonacci numbers in the definition of \(k\) are distinct
and lie strictly between 0 and \(e/2\). Thus no two are congruent
up to sign modulo \(e\), even after multiplication by \(b\).
Their \(r_j\)'s are distinct positive integers with product at least
\(k!\). This proves \eqref{eq:root}.

\section{Cutoffs and exact computation}
It is enough to show
\begin{equation}\label{eq:integer}
  (2F_{n+1})^s k!> 2uvw^2 e^s
\end{equation}
for every possible order. Suppose first that \(m\ge40\).
We have \(F_m^4\le8^m\). Check \(0\le m\le4\) directly. For \(j\ge1\),
\(F_{j+1}\le2F_j\) gives \(F_{j+4}=3F_{j+1}+2F_j\le8F_j\), so
the inequality propagates from \(j\) to \(j+4\). Thus \(uvw^2\le8^m\).

If \(e\le w/32\), then \(2F_{n+1}/e\ge64\).
Since \(r\ge3\), we have \(s=p-\lfloor p/r\rfloor\ge2p/3\).
It follows that \((2F_{n+1}/e)^s k!\ge 2^{4(m-3)}>2\cdot8^m\).
If \(e>w/32\), then \eqref{eq:add} gives
\(w\ge F_{11}F_{m-10}=89F_{m-10}\). Consequently
\(2F_{m-10}<e\), and the indices \(2,\ldots,m-10\) give
\(s\ge k\ge m-11\). Also \(e\le w\le F_{n+1}\).
It follows that \((2F_{n+1}/e)^s k!\ge 2^{m-11}(m-11)!>
2\cdot8^m\). The last strict inequality holds initially since
\(29!>2^{92}\), and its left side gains a factor \(2(m-10)>8\)
in each step. This proves \eqref{eq:integer} for \(m\ge40\).

For \(4\le m<40\) and \(n\ge\max(m,11)\), the left side of
\eqref{eq:integer} is nondecreasing in \(n\). The supplement tests
\(n=\max(m,11)\) and every distinct divisor \(e>1\) of the three
lengths. It obtains all divisors of \(a\) by pairing \(b\) and \(a/b\)
for \(1\le b\le\lfloor\sqrt a\rfloor,\ b\mid a\). The enumeration
uses 1,424 pairs \((m,e)\).

The remaining dimensions are exactly \(4\le m\le n\le10\), giving
28 pairs. Work modulo \(q^{h+1}\) starting from the polynomial~1.
For \(j=1,\ldots,m\), divide by \(1-q^{F_j}\) when \(j>1\)
and multiply by \(1-q^{F_{n+j}}\). The missing denominator
\(1-q^{F_1}\) cancels the leading \(1-q\) in \(\dd P\).
Division by \(1-q^a\) updates coefficients in increasing degree
by \(y_i=x_i+y_{i-a}\). Multiplication uses decreasing degree and
\(y_i=x_i-x_{i-a}\), with zero at negative indices. The surviving
array is exactly \((\dd P(0),\ldots,\dd P(h))\). All its entries
are nonnegative in all 28 cases. This completes the induction and
proves Theorem~\ref{thm:main}.

\section*{Computational supplement}
The accompanying Python 3 program \texttt{original\_verifier.py}
implements the two finite checks in Section~5 using exact integer
arithmetic and the standard library.

\section*{Use of generative AI}
OpenAI's GPT 6.2 Astra contributed to developing the proof and the
computational supplement, checking the argument, and drafting the
manuscript.

\begin{thebibliography}{9}
\bibitem{BCK} N. Bergeron, C. Ceballos and J. K\"ustner.
Elliptic and \(q\)-analogs of the Fibonomial numbers.
\emph{SIGMA} 16 (2020), 076.
\href{https://doi.org/10.3842/SIGMA.2020.076}{doi:10.3842/SIGMA.2020.076}.
\bibitem{CIMSY} B. B. Connelly, E. Ito, T. C. Martinez, O. Shevchenko
and K. Yang. Unimodality of \(q\)-Fibonomial coefficients for small cases.
\href{https://arxiv.org/abs/2605.12822v1}{arXiv:2605.12822v1} (2026).
\end{thebibliography}
\end{document}
